\section{Numerical method}
We advance $\dot y=f(t,y)$ with classical explicit RK4 as presented in Lecture 10 \cite{goza10}. For a step
$h_n=\min(h,T-t_n)$ the update is
\begin{align}
k_1&=f(t_n,y_n), & k_2&=f(t_n+h_n/2,y_n+h_nk_1/2),\\
k_3&=f(t_n+h_n/2,y_n+h_nk_2/2), & k_4&=f(t_n+h_n,y_n+h_nk_3),\\
\widetilde y_{n+1}&=y_n+\frac{h_n}{6}(k_1+2k_2+2k_3+k_4).
\end{align}
Each stage evaluates the derivative at its own time and state, including
forcing and control when the complete model is connected. Four derivative
evaluations are used per step. The last step ends at the specified final time.

To derive the order, augment the nonautonomous system by $\dot t=1$ and
match the elementary differentials in its Taylor expansion. The coefficients are
\[
A=\begin{pmatrix}0&0&0&0\\1/2&0&0&0\\0&1/2&0&0\\0&0&1&0\end{pmatrix},
\quad b=\begin{pmatrix}1/6\\1/3\\1/3\\1/6\end{pmatrix},
\quad c=A\mathbf1=\begin{pmatrix}0\\1/2\\1/2\\1\end{pmatrix}.
\]
Writing powers of $c$ componentwise and $C=\operatorname{diag}(c)$,
the Taylor matching conditions through order four are
\begin{align}
b^T\mathbf1&=1,&b^Tc&=1/2,&b^Tc^2&=1/3,&b^TAc&=1/6,\\
b^Tc^3&=1/4,&b^TCAc&=1/8,&b^TAc^2&=1/12,&b^TA^2c&=1/24.
\end{align}
Direct substitution satisfies all eight conditions. With sufficiently bounded
derivatives through order five and a Lipschitz right hand side in a neighborhood
of the solution, the one step defect from exact initial data is $O(h_n^5)$.
Lecture 11 \cite{goza11} divides this defect by the step length when defining truncation
error. With $\Phi_h$ denoting one unprojected RK4 update, its convention gives
\[
\tau_n=\frac{y(t_n+h_n)-\Phi_{h_n}(t_n,y(t_n))}{h_n}=O(h_n^4).
\]
Thus fifth order one step defect and fourth order truncation error per unit
time describe the same method. The Adams moment conditions in Lecture 11
do not apply directly to RK stage states. The eight conditions above account
for the intermediate states and time dependence.
The error recurrence has the form
\[
\|e_{n+1}\|\leq(1+L_*h_n)\|e_n\|+K h_n^5.
\]
Over a fixed finite interval, summing the propagated defects yields global error
$O(h^4)$. These smoothness assumptions need separate consideration at saturation
crossings. Agreement between refined numerical trajectories alone is not a
verification of the governing equations.

For $\dot y=\lambda y$, substitution of the four stages gives
\[
y_{n+1}=R(z)y_n,\qquad
R(z)=1+z+z^2/2+z^3/6+z^4/24,\qquad z=h\lambda.
\]
Following Lecture 12 \cite{goza12}, strict modal decay requires $|R(z)|<1$.
Equality is the marginal boundary, not strict decay. On the imaginary axis,
$|R(i\xi)|^2=1-\xi^6/72+\xi^8/576$, so the stable interval is
$|\xi|\leq2\sqrt2$. This linear result does not establish nonlinear accuracy.
On the negative real axis the bounded region extends to approximately
$z=-2.785$. Consequently a fast decaying mode can restrict an explicit step
even when its amplitude is small. Lecture 13 \cite{goza13} distinguishes that stability
restriction from the step needed to resolve the desired dynamics.
Frozen Jacobian spectra should be reviewed at representative trajectory
states and across gains. Neutral modes must be counted separately when
forming a ratio of nonzero mode magnitudes. Nonnormality, time dependence,
projection and saturation prevent this check from being a full nonlinear
stability proof. No full model stiffness classification has been measured.
For a frozen field, the unsaturated rotational part of the normalized feedback
has torque $-k_B(I-bb^T)\omega$, where $b=B_B/\|B_B\|$.
The preliminary gain and inertia give a largest damping scale
$k_B/J_{\min}=1.43\times10^{-3}\ \mathrm{s}^{-1}$.
At the preliminary rate of $10$ degrees per second, the quaternion frequency
is $\|\omega\|/2=0.0873\ \mathrm{s}^{-1}$.
These are scale estimates, not a stiffness classification or an accepted step.
Orbital forcing, rigid body coupling, gain sweeps, and saturation still require
the complete model assessment.

After each completed step we record $\|\widetilde q_{n+1}\|-1$, then set
$q_{n+1}=\widetilde q_{n+1}/\|\widetilde q_{n+1}\|$ while leaving angular
velocity unchanged. The input quaternion must already have unit norm.
Projection is smooth near a unit quaternion and fixes the exact solution,
so it preserves fourth order there. Intermediate stages are not projected by
the integrator. The verification fixture uses polynomial quaternion kinematics
and evaluates no attitude matrix in its derivative. The shared dynamics module uses a normalized local copy for rotation matrix
evaluation and retains the unprojected stage value in quaternion kinematics.
The new integration checks exercise that contract with artificial zero torque
control. Independent real field and control checks and full model refinement are reported in the verification section.

\begin{enumerate}
\item Evaluate the four derivatives using the stage times and states.
\item Form the weighted update and record the norm error before projection.
\item Project only the completed quaternion and store the state and time.
\item Repeat until the deadline, shortening the final step when needed.
\end{enumerate}
RK4 offers fourth order accuracy for smooth trajectories with four evaluations
per step and no nonlinear solve. It is a candidate for this model because the
right hand side is small and explicit. Its measured baseline accuracy and cost are reported with the full model refinement; extrapolation to other engineering cases requires separate checks. The present implementation stores
all accepted states, using memory proportional to the number of steps.
